\section*{1.2. Các loại bài toán máy học}

\paragraph{Câu 1.} Phân biệt học có giám sát, học không giám sát và học tăng cường.

\medskip\noindent\textbf{Trả lời.}
\begin{itemize}[nosep]
\item \textbf{Học có giám sát:} dữ liệu gồm nhãn; mô hình học ánh xạ từ đặc trưng đến nhãn. Kiểu đầu ra: nhãn rời rạc (phân lớp) hoặc liên tục (hồi quy).
\item \textbf{Học không giám sát:} dữ liệu không có nhãn; tập trung vào cấu trúc nội tại (phân cụm, giảm chiều, mô hình sinh, phát hiện bất thường theo mật độ, \ldots).
\item \textbf{Học tăng cường:} tác tử tương tác với \textbf{môi trường} qua hành động; nhận tín hiệu thưởng/phạt; mục tiêu học \textbf{chính sách} tối đa hóa phần thưởng tích lũy theo thời gian. Điểm khác có giám sát là tín hiệu thường \textbf{thưa, trễ} và phụ thuộc chuỗi hành động.
\end{itemize}

\paragraph{Câu 2.} Cho ví dụ thực tế cho từng loại bài toán máy học.

\medskip\noindent\textbf{Trả lời.}
\begin{itemize}[nosep]
\item \textbf{Có giám sát:} phát hiện gian lận thẻ khi có nhãn giao dịch ``gian lận/tốt''; dự báo mưa theo ảnh vệ tinh có đo thực địa.
\item \textbf{Không giám sát:} phân nhóm khách hàng theo lịch sử mua hàng khi chưa gắn nhãn nhóm; giảm chiều dữ liệu gene để trực quan hóa.
\item \textbf{Tăng cường:} huấn luyện robot đi bộ trong mô phỏng; chơi cờ với phần thưởng thắng/thua; điều khiển hệ thống tiết kiệm năng lượng với chi phí chạy máy.
\end{itemize}

\paragraph{Câu 3.} Phân biệt bài toán phân lớp (classification) và hồi quy (regression).

\medskip\noindent\textbf{Trả lời.} Cả hai đều thuộc học có giám sát, nhưng \textbf{dạng nhãn} khác nhau:
\begin{itemize}[nosep]
\item \textbf{Phân lớp:} nhãn rời rạc (lớp); mô hình dự đoán lớp hoặc phân phối xác suất trên các lớp.
\item \textbf{Hồi quy:} nhãn (mục tiêu) là số thực (hoặc vector số); mô hình dự đoán giá trị số (ví dụ giá nhà, nồng độ).
\end{itemize}
Ngưỡng hóa đầu ra hồi quy thành khoảng rời rạc đôi khi được dùng trong thực tế, nhưng về bản chất bài toán vẫn được xếp theo loại nhãn gốc.

\paragraph{Câu 4.} Phân biệt clustering và classification.

\medskip\noindent\textbf{Trả lời.}
\begin{itemize}[nosep]
\item \textbf{Phân lớp:} học có giám sát; \textbf{có nhãn lớp} cho tập huấn luyện; mục tiêu gán lớp cho mẫu mới theo nhãn hiện có.
\item \textbf{Phân cụm (clustering):} học không giám sát; \textbf{không có nhãn}; mục tiêu tìm nhóm mẫu tương đồng (cụm) dựa trên khoảng cách hoặc mật độ.
\end{itemize}
Kết quả cụm không nhất thiết trùng với ``lớp nghiệp vụ''; cần diễn giải hoặc đối chiếu thêm với chuyên môn.

\paragraph{Câu 5.} Bài toán dự đoán giá nhà thuộc loại bài toán nào? Vì sao?

\medskip\noindent\textbf{Trả lời.} Đây là bài toán \textbf{hồi quy trong học có giám sát}. Mục tiêu là dự đoán \textbf{giá nhà}---một biến liên tục (hoặc rất nhiều mức giá khác nhau)---từ các đặc trưng đầu vào (diện tích, vị trí, \ldots), khi đã có lịch sử các cặp đặc trưng--giá thực tế để huấn luyện.
